\chapter{المناطق الخالية من الأصفار والأصفار الاستثنائية}

\section{نطاق الفصل وحالته}

\noindent\textbf{حالة الفصل:} \textenglish{REVIEWED}. يطوّر هذا الفصل عدم
الانعدام النوعي على الخط \(\Re(s)=1\)، المثبت في الفصل العاشر، إلى منطقة
كمية قرب الخط لدوال ديريشليه \(L\). كما يعزل الصفر الحقيقي الاستثنائي
الممكن، ويثبت فرادته عبر الموصلات المحدودة بمبرهنة Landau--Page، ثم يفصل
بين مبرهنة Siegel غير الفعالة وظاهرة Deuring--Heilbronn لتنافر بقية
الأصفار.

بدأ التأليف بعد إغلاق بوابة
\textenglish{EVIDENCE-FIRST / PRE-AUTHORING} بحكم \textenglish{PASS} في
\texttt{docs/CHAPTER\_11\_LOGIC\_AUDIT\_2026-07-20.md}، ثم اجتاز تدقيق
ما بعد التأليف بحكم \textenglish{PASS} في
\texttt{docs/CHAPTER\_11\_AUTHORING\_AUDIT\_2026-07-20.md}. ولا يدعي
الفصل أفضل ثابت عددي، ولا يثبت البرهان الكامل لظاهرة
Deuring--Heilbronn، ولا يتناول Siegel--Walfisz أو Bombieri--Vinogradov أو
Linnik، ولا يستعمل فرضية ريمان المعممة. وتبقى منطقة Vinogradov--Korobov
وتقديرات كثافة الأصفار موضوعين متقدمين مؤجلين.

\subsection{المتطلبات السابقة}

يعتمد الفصل على:
\begin{itemize}
  \item الشخصيات البدائية والموصل والدالة المكتملة والمعادلة الوظيفية من
  الفصل السابع.
  \item صيغة ستيرلينغ المنتظمة من الفصل السادس.
  \item المتراجحة الموزونة وعدم الانعدام على الخط \(\Re(s)=1\) من الفصل
  العاشر.
  \item جداء هادامار للدوال التامة من الرتبة الأولى بوصفه نتيجة قياسية في
  التحليل المركب.
\end{itemize}

\subsection{نواتج التعلم}

بعد إتمام الفصل ينبغي أن يستطيع القارئ:
\begin{enumerate}
  \item تفسير ظهور المقياس \(\log(q(|t|+2))\) قرب الخط الواحد.
  \item اشتقاق صيغة المشتقة اللوغاريتمية بمجموع على الأصفار.
  \item إثبات المنطقة القياسية مع استثناء حقيقي بسيط ممكن.
  \item التمييز بين الفرادة داخل دالة واحدة، والفرادة عند ترديد ثابت،
  وفرادة Landau--Page عبر الموصلات.
  \item شرح مبرهنة Siegel ومعنى عدم فعالية ثابتها.
  \item تفسير ظاهرة Deuring--Heilbronn بوصفها تنافرًا كميًا للأصفار.
\end{enumerate}

\section{المقياس التحليلي وثلاثة معاني للفرادة}

لتكن \(\chi\) شخصية بدائية غير رئيسية بموصل \(q\ge3\)، ولتكن
\[
s=\sigma+it.
\]
نستعمل المقياس
\[
\mathcal L_q(t)=\log\bigl(q(|t|+2)\bigr).
\]
يجمع هذا المقياس كلفة الموصل \(\log q\) وكلفة عامل غاما عند الارتفاع
\(|t|\). ويجب فصل العبارات الآتية:
\begin{enumerate}
  \item داخل دالة بدائية واحدة، يوجد على الأكثر صفر حقيقي واحد في المنطقة
  القياسية، ويكون بسيطًا.
  \item عند ترديد ثابت، يملك حاصل ضرب جميع دوال الشخصيات على الأكثر صفرًا
  استثنائيًا واحدًا في منطقة صريحة مناسبة؛ انظر
  \cite{McCurley1984,Kadiri2018}.
  \item عبر الموصلات حتى \(Q\)، تمنع مبرهنة Landau--Page وجود شخصيتين
  بدائيتين مختلفتين لهما صفران شديدا القرب من الواحد.
\end{enumerate}

\section{الرد إلى الشخصية البدائية}

\begin{proposition}[الأصفار القريبة من الواحد والجد البدائي]
\label{prop:chapter11-primitive-reduction}
\resultid{ANT-PROP-11-01}
\provedhere
لتكن \(\chi\pmod q\) شخصية مستحثة من شخصية بدائية
\(\chi^*\pmod r\). عندئذ
\[
L(s,\chi)=L(s,\chi^*)
\prod_{\substack{p\mid q\\p\nmid r}}
\left(1-\chi^*(p)p^{-s}\right).
\]
ولا ينعدم أي عامل محلي في \(\Re(s)>0\). كما أن الفرق بين المشتقتين
اللوغاريتميتين هو \(O(\log q)\) عندما \(\Re(s)>1\).
\end{proposition}

\begin{proof}
صيغة العوامل المحلية هي
\textenglish{\texttt{ANT-PROP-07-03}}. وإذا \(\sigma>0\)، فإن
\[
|\chi^*(p)p^{-s}|\le p^{-\sigma}<1,
\]
فلا ينعدم العامل. وبالتفاضل اللوغاريتمي:
\[
\frac{L'}{L}(s,\chi)=\frac{L'}{L}(s,\chi^*)+
\sum_{\substack{p\mid q\\p\nmid r}}
\frac{\chi^*(p)(\log p)p^{-s}}
     {1-\chi^*(p)p^{-s}}.
\]
ولـ\(\sigma>1\):
\[
\left|
\frac{\chi^*(p)(\log p)p^{-s}}
     {1-\chi^*(p)p^{-s}}
\right|
\le \frac{\log p}{p^\sigma-1}\le\log p.
\]
ومن ثم لا يتجاوز مجموع الحدود
\[
\sum_{p\mid q}\log p=\log\operatorname{rad}(q)\le\log q.
\]
\end{proof}

\section{الدالة المكتملة والمشتقة اللوغاريتمية}

ضع
\[
a_\chi=\frac{1-\chi(-1)}2\in\{0,1\},
\]
وعرف
\[
\Lambda(s,\chi)=
\left(\frac q\pi\right)^{(s+a_\chi)/2}
\Gamma\!\left(\frac{s+a_\chi}{2}\right)L(s,\chi).
\]
الدالة \(\Lambda(s,\chi)\) تامة من الرتبة الأولى، وتحقق
\[
\Lambda(s,\chi)=\varepsilon_\chi
\Lambda(1-s,\overline\chi),
\qquad |\varepsilon_\chi|=1.
\]

\begin{lemma}[صيغة المشتقة اللوغاريتمية بمجموع على الأصفار]
\label{lem:chapter11-log-derivative-zeros}
\resultid{ANT-LEM-11-01}
\provedhere
إذا كانت \(\chi\) بدائية وغير رئيسية، و\(\Re(s)>1\)، فإن
\[
-\Re\frac{L'}{L}(s,\chi)=
\frac12\log\frac q\pi+
\frac12\Re\frac{\Gamma'}{\Gamma}
\!\left(\frac{s+a_\chi}{2}\right)
-
\sum_\rho^*\Re\frac1{s-\rho},
\]
حيث يجري الجمع المتماثل على الأصفار غير البديهية. وبخاصة، مساهمة كل صفر
في \(-\Re L'/L\) سالبة.
\end{lemma}

\begin{proof}
تعطي مبرهنة هادامار العامة
\[
\Lambda(s,\chi)=
\exp(A_\chi+B_\chi s)
\prod_\rho
\left(1-\frac{s}{\rho}\right)e^{s/\rho}.
\]
وبالتفاضل اللوغاريتمي والتجميع المتماثل:
\[
\frac{\Lambda'}{\Lambda}(s,\chi)=
 b_\chi+\sum_\rho^*\frac1{s-\rho}.
\]
تجعل المعادلة الوظيفية وتناظر الأصفار الجزء الحقيقي للثابت المنتظم
\(b_\chi\) مساويًا للصفر. ومن تعريف \(\Lambda\):
\[
\frac{\Lambda'}{\Lambda}(s,\chi)=
\frac12\log\frac q\pi+
\frac12\frac{\Gamma'}{\Gamma}
\!\left(\frac{s+a_\chi}{2}\right)+
\frac{L'}{L}(s,\chi).
\]
أخذ الجزء الحقيقي وإعادة الترتيب يعطي الهوية. وإذا
\(s=\sigma+it\) و\(\rho=\beta+i\gamma\)، فإن
\[
\Re\frac1{s-\rho}=
\frac{\sigma-\beta}
{(\sigma-\beta)^2+(t-\gamma)^2}>0.
\]
\end{proof}

\begin{lemma}[ضبط الباقي المنتظم]
\label{lem:chapter11-uniform-remainder}
\provedhere
يوجد ثابت مطلق \(C>0\) بحيث لكل شخصية بدائية غير رئيسية \(\chi\)، ولكل
\(1<\sigma\le2\):
\[
-\Re\frac{L'}{L}(\sigma+it,\chi)
\le C\mathcal L_q(t).
\]
وإذا عزلنا أصفارًا محددة، أمكن إبقاء مساهماتها السالبة صراحة. أما الشخصية
الرئيسية \(\chi_0\pmod q\)، فتخضع للحد
\[
-\Re\frac{L'}{L}(\sigma+it,\chi_0)
\le \Re\frac1{\sigma-1+it}+C\mathcal L_q(t).
\]
\end{lemma}

\begin{proof}
تعطي صيغة ستيرلينغ المنتظمة، في الشريط \(1\le\sigma\le2\)،
\[
\left|
\Re\frac{\Gamma'}{\Gamma}
\!\left(\frac{\sigma+a_\chi+it}{2}\right)
\right|\ll\log(|t|+2).
\]
وحذف مجموع الأصفار المطروح من اللمّة السابقة يعطي حدًا علويًا من رتبة
\(\log q+\log(|t|+2)\). وللشخصية الرئيسية نكتب
\[
L(s,\chi_0)=\zeta(s)\prod_{p\mid q}(1-p^{-s}),
\]
فنحتفظ بقطب زيتا صراحة، ونضبط غاما والعوامل المحلية كما في القضية السابقة.
\end{proof}

\section{المتراجحة الموزونة}

نستعمل النتيجة \textenglish{\texttt{ANT-LEM-10-01}}:
\[
-3\frac{\zeta'}{\zeta}(\sigma)
-4\Re\frac{L'}{L}(\sigma+it,\chi)
-\Re\frac{L'}{L}(\sigma+2it,\chi^2)
\ge0
\qquad(\sigma>1).
\tag{11.1}
\]
مصدرها الموجبية
\[
3+4\Re z+\Re(z^2)\ge0
\qquad(|z|\le1).
\]

\section{المنطقة القياسية الخالية}

\begin{theorem}[المنطقة القياسية والاستثناء الحقيقي]
\label{thm:chapter11-standard-zero-free-region}
\resultid{ANT-THM-11-01}
\provedhere
يوجد ثابت مطلق \(c_0>0\) بحيث لكل شخصية ديريشليه بدائية غير رئيسية
\(\chi\) بموصل \(q\ge3\)، تملك \(L(s,\chi)\) على الأكثر صفرًا واحدًا في
المنطقة
\[
\Re(s)>
1-\frac{c_0}{\log(q(|\Im(s)|+2))}.
\tag{11.2}
\]
إذا وجد الصفر، فهو حقيقي وبسيط، وتكون \(\chi\) شخصية حقيقية تربيعية.
\end{theorem}

\begin{proof}
نكتب
\[
\delta=\sigma-1>0,
\qquad \eta=1-\beta\ge0,
\qquad \mathcal L=\log(q(|\gamma|+2)),
\]
ونختار \(\delta=A/\mathcal L\)، حيث \(A>0\) ثابت مطلق صغير.

\paragraph{الشخصية غير الحقيقية.}
إذا كان \(\rho=\beta+i\gamma\) صفرًا وكانت \(\chi\) غير حقيقية، فإن
\(\chi^2\) غير رئيسية. نأخذ \(t=\gamma\) في (11.1)، ونعزل الصفر، فنحصل
على
\[
0\le
\frac3\delta-
\frac4{\delta+\eta}+C\mathcal L.
\tag{11.3}
\]
إذا افترضنا \(\eta\le c/\mathcal L\)، صار الطرف الأيمن بعد القسمة على
\(\mathcal L\) لا يزيد على
\[
\frac3A-\frac4{A+c}+C.
\]
نختار \(A\) بحيث \(1/A>2C\)، ثم نختار \(c>0\) صغيرًا بما يكفي، فيصبح
التعبير سالبًا، وهو تناقض.

\paragraph{الشخصية الحقيقية والصفر غير الحقيقي.}
إذا كانت \(\chi\) حقيقية و\(\gamma\ne0\)، فإن
\(\overline\rho=\beta-i\gamma\) صفر أيضًا، و\(\chi^2\) رئيسية. بعد عزل
الصفرين وقطب الحد الثالث:
\[
0\le
\frac3\delta+
\frac{\delta}{\delta^2+4\gamma^2}-
\frac4{\delta+\eta}-
\frac{4(\delta+\eta)}{(\delta+\eta)^2+4\gamma^2}
+C\mathcal L.
\tag{11.4}
\]
ضع
\[
r=\frac{\delta+\eta}{\delta},
\qquad v=\frac{2|\gamma|}{\delta}.
\]
بعد ضرب الجزء الكسري في \(\delta\)، يصبح
\[
G(r,v)=3+\frac1{1+v^2}-\frac4r-
\frac{4r}{r^2+v^2}.
\]
لدينا
\[
G(1,v)=-1-\frac3{1+v^2}\le-1,
\qquad
\left|\frac{\partial G}{\partial r}(r,v)\right|\le8.
\]
فإذا \(\eta\le\delta/16\)، أعطت مبرهنة القيمة المتوسطة
\(G(r,v)\le-1/2\). ومن ثم
\[
0\le-\frac1{2\delta}+C\mathcal L,
\]
وهو تناقض عندما \(1/(2A)>C\).

\paragraph{الصفر الحقيقي وبساطته.}
افترض وجود أصفار حقيقية \(\beta_j=1-\eta_j\) في المنطقة، محسوبة مع
الرتب. عند \(t=0\)، يعطي قطبا الحدين الأول والثالث \(4/\delta\)، وتعطي
الأصفار
\[
-4\sum_j\frac1{\delta+\eta_j}.
\]
إذا كان مجموع الرتب على الأقل اثنين، وكانت \(\eta_j\le\delta/2\)، فإن
\[
-4\sum_j\frac1{\delta+\eta_j}
\le-\frac{16}{3\delta}.
\]
فتصبح المتراجحة
\[
0\le-\frac4{3\delta}+C\log(2q),
\]
وتناقض اختيار \(A\) صغيرًا. إذن يوجد على الأكثر صفر حقيقي واحد، وتكون
رتبته واحدًا. وبجمع الثوابت الصغيرة المستعملة نحصل على \(c_0\).
\end{proof}

\begin{remark}[الثوابت الصريحة]
أثبت McCurley منطقة صريحة لحاصل ضرب دوال الشخصيات بترديد ثابت
\cite{McCurley1984}، وحسن Kadiri ثوابت عددية في نطاقات محددة
\cite{Kadiri2018}. لا يجوز نقل ثابت من صيغة إلى تطبيع لوغاريتمي مختلف من
دون إعادة الحساب.
\end{remark}

\section{مبرهنة Landau--Page}

\begin{theorem}[مبرهنة Landau--Page النوعية]
\label{thm:chapter11-landau-page}
\resultid{ANT-THM-11-02}
\provedhere
يوجد ثابت مطلق \(c_P>0\) بحيث لكل \(Q\ge3\)، توجد على الأكثر شخصية
ديريشليه حقيقية بدائية واحدة \(\chi\) موصلها \(q\le Q\) وتملك صفرًا
حقيقيًا \(\beta\) يحقق
\[
1-\frac{c_P}{\log Q}<\beta<1.
\]
\end{theorem}

\begin{proof}
نفترض وجود شخصيتين حقيقيتين بدائيتين متميزتين
\(\chi_i\pmod{q_i}\)، حيث \(q_i\le Q\)، ولكل منهما صفر
\(\beta_i=1-\eta_i\) مع \(\eta_i\le c/\log Q\). على الترديد
\(m=\operatorname{lcm}(q_1,q_2)\)، ضع
\[
\chi(n)=\chi_1(n)\chi_2(n).
\]
هي غير رئيسية، وإلا اتفق جدا الشخصيتين البدائيان. نعرف
\[
F(s)=\zeta(s)L(s,\chi_1)L(s,\chi_2)L(s,\chi).
\]
في \(\sigma>1\):
\[
-\frac{F'}{F}(\sigma)=
\sum_{n\ge1}\frac{\Lambda(n)}{n^\sigma}
(1+\chi_1(n))(1+\chi_2(n))\ge0.
\]
نحتفظ بشخصية حاصل الضرب كما هي للمحافظة على الموجبية. وعند تقدير مشتقتها
نردها إلى جدها البدائي، الذي لا يتجاوز موصله \(q_1q_2\le Q^2\).

بوضع \(\delta=\sigma-1\) وعزل الصفرين نحصل على
\[
0\le
\frac1\delta-
\frac1{\delta+\eta_1}-
\frac1{\delta+\eta_2}+C\log Q.
\tag{11.5}
\]
نختار \(\delta=A/\log Q\). وإذا كانت
\(\eta_i\le c/\log Q\)، صار الطرف الأيمن بعد القسمة على \(\log Q\) لا
يزيد على
\[
\frac1A-\frac2{A+c}+C.
\]
نختار \(A\) بحيث \(1/A>2C\)، ثم \(c_P\) صغيرًا بما يكفي، فنحصل على
تناقض.
\end{proof}

\begin{remark}[ما الذي تثبته Page؟]
لا تثبت المبرهنة وجود صفر استثنائي، بل فرادته إذا ظهر. انظر
\cite{Davenport2000,BasakPratt2026}.
\end{remark}

\section{مبرهنة Siegel وعدم الفعالية}

\begin{theorem}[مبرهنة Siegel]
\label{thm:chapter11-siegel}
\resultid{ANT-THM-11-03}
\citedresult{\cite{Siegel1935,Liu2022Siegel}}
لكل \(\varepsilon>0\)، يوجد ثابت \(c(\varepsilon)>0\) بحيث لكل شخصية
ديريشليه حقيقية بدائية غير رئيسية \(\chi\) بموصل \(q\):
\[
|L(1,\chi)|\ge c(\varepsilon)q^{-\varepsilon}.
\tag{11.6}
\]
الثابت \(c(\varepsilon)\) غير فعال.
\end{theorem}

\begin{lemma}[حد المشتقة قرب الواحد]
\label{lem:chapter11-derivative-near-one}
\provedhere
يوجد ثابت مطلق \(C>0\) بحيث، إذا كانت \(\chi\) غير رئيسية بترديد \(q\)،
وكان
\[
1-\frac{c_0}{\log(2q)}\le\sigma\le1,
\]
فإن
\[
|L'(\sigma,\chi)|\le C\log^2(2q).
\]
\end{lemma}

\begin{proof}
ضع \(S(x)=\sum_{n\le x}\chi(n)\). من الدورية وانعدام مجموع دورة كاملة
لدينا \(|S(x)|\le q\). وبالجمع الجزئي عند \(q\):
\[
L(s,\chi)=
\sum_{n\le q}\frac{\chi(n)}{n^s}+
 s\int_q^\infty\frac{S(x)}{x^{s+1}}\,dx
\qquad(\Re(s)>0).
\]
وبالتفاضل عند \(s=\sigma\) الحقيقي:
\begin{align*}
L'(\sigma,\chi)={}&-
\sum_{n\le q}\frac{\chi(n)\log n}{n^\sigma}+
\int_q^\infty\frac{S(x)}{x^{\sigma+1}}\,dx\\
&-
\sigma\int_q^\infty
\frac{S(x)\log x}{x^{\sigma+1}}\,dx.
\end{align*}
في المجال المحدد \(q^{1-\sigma}\le e^{c_0}\). لذلك لا يتجاوز الحد
المنتهي \(O(\log^2(2q))\)، ولا يتجاوز التكاملان تباعًا
\(O(1)\) و\(O(\log(2q))\).
\end{proof}

\begin{corollary}[الإبعاد غير الفعال للصفر الاستثنائي]
\label{cor:chapter11-siegel-zero-distance}
\resultid{ANT-COR-11-01}
\provedhere
لكل \(\varepsilon>0\)، يوجد ثابت غير فعال \(c_\varepsilon>0\) بحيث إذا
كان \(\beta\) صفرًا استثنائيًا لدالة شخصية حقيقية بدائية بموصل \(q\)، فإن
\[
1-\beta\ge c_\varepsilon q^{-\varepsilon}.
\]
\end{corollary}

\begin{proof}
لدينا
\[
L(1,\chi)=\int_\beta^1L'(u,\chi)\,du,
\]
فتعطي اللمّة السابقة
\[
|L(1,\chi)|\ll(1-\beta)\log^2(2q).
\]
نطبق مبرهنة Siegel بالأس \(\varepsilon/2\)، ثم نستعمل
\(\log^2(2q)\ll_\varepsilon q^{\varepsilon/2}\)، فنحصل على النتيجة. ويرث
الثابت عدم الفعالية من مبرهنة Siegel.
\end{proof}

\begin{remark}[منع إساءة استعمال مبرهنة Siegel]
كل نتيجة تعتمد على ثابت Siegel من دون بديل فعال ترث عدم الفعالية. لذلك لا
يجوز إدخال هذا الثابت في حد خطأ ثم وصف الحد بأنه فعال.
\end{remark}

\section{ظاهرة Deuring--Heilbronn}

للترديد الثابت \(q\)، ضع
\[
\mathcal D_q(s)=\prod_{\chi\bmod q}L(s,\chi).
\]

\begin{theorem}[ظاهرة Deuring--Heilbronn، صيغة نوعية]
\label{thm:chapter11-deuring-heilbronn}
\resultid{ANT-THM-11-04}
\citedresult{\cite{BenliGoelTwissZaman2026}}
توجد ثوابت مطلقة فعالة موجبة \(c_0,c_1,c_2\) بحيث إذا كان \(T\ge2\)،
وكانت \(\mathcal D_q\) تملك صفرًا استثنائيًا \(\beta_1\) يحقق
\[
\lambda_1=(1-\beta_1)\log(q(T+2))\le c_0,
\]
فإن كل صفر آخر \(\rho=\beta+i\gamma\) لـ\(\mathcal D_q\)، مع
\(|\gamma|\le T\)، يحقق
\[
1-\beta\ge
\frac{c_1}{\log(q(T+2))}
\log\!\left(\frac{c_2}{\lambda_1}\right).
\tag{11.7}
\]
\end{theorem}

\noindent لا ندعي هنا البرهان الكامل، الذي يستعمل أدوات تتجاوز الفصول
الحالية.

\begin{remark}[الفرادة ليست التنافر]
Landau--Page تمنع استثناءين مستقلين من الاقتراب معًا من الواحد. أما
Deuring--Heilbronn فتدفع كل صفر آخر بعيدًا بمقدار يزداد عندما يقترب
\(\beta_1\) أكثر من الواحد.
\end{remark}

\section{الأثر في توزيع الأوليات}

تشرح الصيغة الصريحة الإرشادية
\[
\psi(x,\chi)\approx-
\sum_\rho\frac{x^\rho}{\rho}+\text{حدود أخرى}
\]
سبب خطورة الصفر الاستثنائي. فإذا كان \(\beta\) قريبًا من الواحد، ظهر حد
كبير من رتبة \(-x^\beta/\beta\).

\begin{proposition}[خريطة الأثر في الصيغة الصريحة]
\label{prop:chapter11-explicit-formula-diagnostic}
\resultid{ANT-PROP-11-02}
\deferredresult{الصيغة الصريحة الكاملة لدوال الشخصيات مؤجلة إلى فصل كمي لاحق}
تضبط المنطقة القياسية مساهمة الأصفار غير الاستثنائية، بينما يساهم الصفر
الحقيقي الاستثنائي بحد منفرد يجب إظهاره صراحة. وتمنح ظاهرة
Deuring--Heilbronn تعويضًا جزئيًا بدفع بقية الأصفار بعيدًا.
\end{proposition}

\section{حدود الفعالية}

\begin{center}
\begin{tabular}{|p{0.28\textwidth}|p{0.19\textwidth}|p{0.38\textwidth}|}
\hline
\textbf{النتيجة} & \textbf{الفعالية} & \textbf{الملاحظة} \\
\hline
المنطقة القياسية الفردية & فعالة من حيث المبدأ & الثابت غير محسن \\
\hline
Landau--Page النوعية & فعالة من حيث المبدأ & تثبت الفرادة فقط \\
\hline
مبرهنة Siegel & غير فعالة & ثابت \(c(\varepsilon)\) غير معلوم حسابيًا \\
\hline
الإبعاد \(1-\beta\gg_\varepsilon q^{-\varepsilon}\) & غير فعال & يرث ثابت Siegel \\
\hline
Deuring--Heilbronn الصريحة & فعالة & مقتبسة؛ البرهان الكامل مؤجل \\
\hline
\end{tabular}
\end{center}

\section{تمارين}

\begin{enumerate}
  \item أثبت أن
  \[
  \sum_{p\mid q}\frac{\log p}{p^\sigma-1}\le\log q
  \qquad(\sigma>1).
  \]
  \item تحقق من أن الشخصية غير الحقيقية لا يمكن أن يكون مربعها شخصية
  رئيسية.
  \item أثبت الحد \(|\partial G/\partial r|\le8\) المستعمل في برهان
  المنطقة القياسية.
  \item استخرج قيمًا صريحة غير مثلى للثوابت بدلالة ثابت الباقي \(C\).
  \item أثبت أن شخصية حاصل الضرب في برهان Landau--Page غير رئيسية عندما
  تكون \(\chi_1,\chi_2\) بدائيتين متميزتين.
  \item تحقق من الموجبية
  \[
  1+\chi_1(n)+\chi_2(n)+\chi_1(n)\chi_2(n)
  =(1+\chi_1(n))(1+\chi_2(n))\ge0.
  \]
  \item أثبت حد المشتقة قرب الواحد باستعمال دورية الشخصية والجمع الجزئي.
  \item اشرح لماذا لا تنتج مبرهنة Siegel خوارزمية لحصر جميع الأصفار
  الاستثنائية حتى موصل معلوم.
  \item قارن بين مستويات الفرادة الثلاثة المذكورة في بداية الفصل.
  \item اشرح من الصيغة الإرشادية لـ\(\psi(x,\chi)\) خطورة الحد
  \(x^\beta/\beta\).
\end{enumerate}

\section{خلاصة الفصل}

حوّلت المشتقة اللوغاريتمية هندسة الأصفار إلى متراجحات حقيقية: كل صفر يعطي
مساهمة سالبة، بينما لا تتجاوز كلفة الموصل وغاما رتبة
\(\log(q(|t|+2))\). وبضم ذلك إلى كثيرة حدود مثلثية غير سالبة حصلنا على
المنطقة القياسية، مع إمكان استثناء حقيقي بسيط لشخصية تربيعية.

ثم أثبتنا مبرهنة Landau--Page من حاصل ضرب أويلري موجب، وفصلنا مبرهنة
Siegel غير الفعالة عن النتائج الفعالة، وسجلنا ظاهرة Deuring--Heilbronn
بوصفها نتيجة مقتبسة لتنافر بقية الأصفار. وهذه الأدوات هي المدخل الطبيعي
لمبرهنة Siegel--Walfisz وتقديرات الكثافة وBombieri--Vinogradov.